%%%%

%%%   Самарский государственный технический университет

%%%   Кафедра прикладной математики и информатики

%%%

%%%   Преамбула для статьи в журнал "Вестник Самарского государственного

%%%   технического университета. Серия: «Физико-математические науки»"

%%%   Ориентирована под MikTEx 2.4 for Win32

%%%

%%%   Хотя сам журнал собирается под GNU/Linux на дистрибутиве TeXLive



\documentclass[11pt,a4paper,twoside]{article}



\usepackage[cp1251]{inputenc}

\usepackage[T2A]{fontenc}

\usepackage[english,russian]{babel}

\usepackage{samgtubib}

\usepackage[dvips]{graphicx}

\usepackage[multiple]{footmisc}



\usepackage{samgtu}

\renewcommand{\VestnikYear}{2009} % Год издания Вестника

\renewcommand{\VestnikNumber}{2\,(19)} % Номер Вестника

\renewcommand{\VestnikMonth}{Ноябрь} % Месяц выхода Вестника

\renewcommand{\VestnikData}{01.11.09} % Дата подписания Вестника в печать

\renewcommand{\VestnikUPL}{26,64} % Усл. печ. л.

\renewcommand{\VestnikUIL}{25,73} % Уч.-изд. л.

\renewcommand{\VestnikRegNumber}{479/09} % Регистрационный номер

\renewcommand{\VestnikOrder}{994} % Номер заказа







%

%\begin{document}





\renewcommand{\firstpage}{82}

\renewcommand{\lastpage}{89}



\udk{517.957}

\msc{35K35, 35K61}



\title{Решения анизотропных параболических уравнений с~двойной нелинейностью в~неограниченных областях}

{Решения анизотропных параболических

уравнений \dots}



\author{Л.\,М.~Кожевникова, А.\,А.~Леонтьев}

{Л.~М.~К\,о\,ж\,е\,в\,н\,и\,к\,о\,в\,а, А.~А.~Л\,е\,о\,н\,т\,ь\,е\,в }



\institute{\bashgusterl.

}



\email{E-mails}{kosul@mail.ru, alexey\_leontiev@inbox.ru}



\otherinf{\fullauthor{Лариса Михайловна Кожевникова}\regalia{д.ф.-м.н., проф.} \position{профессор}

\dept{каф. математического анализа}

\fullauthor{Алексей Александрович Леонтьев} \position{аспирант} \dept{каф. математического анализа}}



\keywords{анизотропное уравнение,

параболическое уравнение с двойной нелинейностью, существование

решения, скорость убывания решения}



\workabstract{

 Работа посвящена некоторому классу

анизотропных параболических уравнений высокого порядка с двойной нелинейностью,

представителем которого является модельное уравнение вида

\begin{gather*}

\frac{\partial}{\partial t}\left(|u|^{k-2}u \right)=

\sum_{\alpha=1}^n(-1)^{m_\alpha-1}\frac{\partial^{m_\alpha}}{\partial x_\alpha^{m_\alpha}}

\left[\left|\frac{\partial^{m_\alpha} u}{\partial x_\alpha^{m_\alpha}}\right|^{p_\alpha-2}

\frac{\partial^{m_\alpha} u}{\partial x_\alpha^{m_\alpha}}\right],\\

m_1,\ldots, m_n\in \mathbb{N},\quad

p_n\geq \ldots \geq p_1>k,\quad k>1.

\end{gather*}

Для решений первой

смешанной задачи в цилиндрической области

$D=(0,\infty)\times\Omega$ с неограниченной областью $\Omega\subset \mathbb{R}_n$, $n\geq 2$,

с однородным краевым условием Дирихле и финитной начальной

функцией установлена максимальная скорость  убывания  при

$t \to \infty$. Ранее авторами были получены оценки сверху

 для анизотропных уравнений второго порядка и доказана их точность.}



\engtitle{Solutions of anisotropic parabolic equations with double non-linearity in unbounded domains}



\engauthor{L.\,M.~Kozhevnikova, A.\,A.~Leontiev}{L.~M.~K\,o\,z\,h\,e\,v\,n\,i\,k\,o\,v\,a,

A.~A.~L\,e\,o\,n\,t\,i\,e\,v}





\enginstitute{\engbashgusterl.

}



\otherenginfm{\fullauthor{Larisa M. Kozhevnikova} \regalia{Dr.\,Sci. (Phys. \& Math.)}

\dept{Dept. of Mathematical Analysis} \\

\fullauthor{Alexey~A. Leontiev}

\position{Postgraduate Student} \dept{Dept. of Mathematical Analysis}}



\engabstract{This work is devoted to some class of parabolic equations of high order with double nonlinearity which can be represented by a model equation \begin{gather*} \frac{\partial}{\partial t}(|u|^{k-2}u)= \sum_{\alpha=1}^n(-1)^{m_\alpha-1}\frac{\partial^{m_\alpha}}{\partial x_\alpha^{m_\alpha}} \left[\left|\frac{\partial^{m_\alpha} u}{\partial x_\alpha^{m_\alpha}}\right|^{p_\alpha-2} \frac{\partial^{m_\alpha} u}{\partial x_\alpha^{m_\alpha}}\right],\\ m_1,\ldots, m_n\in \mathbb{N},\quad p_n\geq \ldots \geq p_1>k,\quad k>1. \end{gather*} For the solution of the first mixed problem in a cylindrical domain $ D=(0,\infty)$ $\times\Omega, \;\Omega\subset \mathbb{R}_n,$ $n\geq 2,$ with homogeneous Dirichlet boundary condition and finite initial function the highest rate of decay established as $t \to \infty$. Earlier upper estimates were obtained by the authors for anisotropic equation of the second order and prove their accuracy.}



\engkeywords{anisotropic equation, doubly nonlinear

parabolic equations, existence of strong solution, decay rate of

solution}



\date{15/XI/2012}

\revisiondate{10/III/2013}









\maketitle



\smallskip



\Section[n]{Введение}  Пусть $\Omega $ "--- неограниченная область пространства

 $\mathbb{R}_{n}=\{{\bf x}=(x_{1},x_{2},\ldots,x_{n})\}$, $n \geq 2$.

 В цилиндрической области $D=\{t>0\}\times \Omega$ для анизотропного квазилинейного параболического уравнения высокого порядка

 рассматривается первая смешанная задача

\begin{equation}

\label{eq:kojleont:ur}

\bigl(|u|^{k-2}u \bigr)_t=

\sum_{\alpha=1}^n(-1)^{m_\alpha-1}\frac{\partial^{m_\alpha}}{\partial x_\alpha^{m_\alpha}}

\left[a_{\alpha}\left(\Bigl(\frac{\partial^{m_\alpha} u}{\partial x_\alpha^{m_\alpha}}\Bigr)^2\right)

\frac{\partial^{m_\alpha} u}{\partial x_\alpha^{m_\alpha}}\right],\quad(t,{\bf x}) \in D,

\end{equation}

где $k>1$, $m_\alpha\in \mathbb{N}$;

\begin{gather}

\label{eq:kojleont:gu} D^{j}_{x_\alpha}u(t,{\bf x})\big|_{S}=

0;\quad j= 1, 2, \dots, m_\alpha-1, \; S=\{t>0\}\times\partial \Omega;\\

\label{eq:kojleont:gu} u(0,{\bf x})=\varphi({\bf x}),\quad

\varphi({\bf x})\in L_k(\Omega),\quad\varphi_{x_{\alpha}}({\bf

x})\in L_{p_{\alpha}}(\Omega).

\end{gather}

Предполагается, что неотрицательные функции $a_{\alpha}(s)$,

$s\geq 0$, $\alpha=1, 2, \dots, n,$  подчиняются следующим условиям:

$a_{\alpha}(0)=0$, $a_{\alpha}(s)\in C^1(0,\infty),$

\begin{gather}

\label{eq:kojleont:us1} \overline{a}s^{(p_{\alpha}-2)/2}\leq

a_{\alpha}(s)\leq\widehat{a} s^{(p_{\alpha}-2)/2},\\

\frac{p_1}{2}a_{\alpha}(s)\leq a_{\alpha}(s)+a'_{\alpha}(s)s\leq

\widehat{b}a_{\alpha}(s),

\end{gather}

с положительными константами $\widehat{a}\geq \overline{a}$,

$2\widehat{b}\geq p_1>k$ ($p_1\leq p_2\leq \ldots\leq p_n$).

Например,

$a_{\alpha}(s)=s^{(p_{\alpha}-2)/2}$, $ \widehat{b}=p_n/2.$



 Исследованию  поведения решений смешанных задач для линейных и квазилинейных параболических уравнений второго и высокого порядков при $t\to \infty$ посвящены работы А.~К.~Гущина, В.~И.~Ушакова, Ф.~Х.~Мукминова,  А.~Ф.~Тедеева, Л.~М.~Кожевниковой, Р.~X.~Каримова, В.~Ф.~Гилимшиной и~др.

Обзоры соответствующих результатов можно найти в~[\citen{kojleont:bib:kojevnik2000}--\citen{kojleont:bib:kar_koj}].



В настоящей работе исследуется допустимая скорость стабилизации решения задачи \eqref{eq:kojleont:ur}--\eqref{eq:kojleont:gu}.

Будем рассматривать области, расположенные вдоль выделенной оси $Ox_s$, $s\in\{1, 2, \dots, n\}$ 

(область $\Omega$ лежит  в~полупространстве $\mathbb{R}^+_n[s]=\{{\bf x}\in \mathbb{R}_n \mid x_s>0\}$, сечение

$\gamma_r[s]=\{\textbf{x}\in\Omega \mid x_s=r\}$ не пусто и ограничено при любом $r>0$). 

Ниже будет использовано следующее обозначение:

$$\Omega_{a}^{b}[s]=\{{\bf x}\in\Omega\;|\;a<x_s<b\},$$ при этом

значения $a=0,\;b=\infty$ опускаются.



Предполагается, что начальная функция имеет ограниченный носитель, так что

\begin{equation}

\label{eq:kojleont:supp}

\mathop{\rm supp} \varphi\subset\Omega^{R_0},\quad R_0>0.

\end{equation}



\smallskip



\begin{theorem}[1]\label{eq:kojleont:the_1(00)}

Пусть выполнено условие  \eqref{eq:kojleont:supp}, тогда найдутся

положительные числа $\kappa,$

$\mathcal{M}(p_s,m_s,k,\overline{a},\widehat{a})$

 такие, что ограниченное решение $u(t,\mathbf{x})$

 задачи \eqref{eq:kojleont:ur}--\eqref{eq:kojleont:gu} при всех $t\geq

0$, $r\geq 2R_0$ удовлетворяет оценке

\begin{equation}\label{eq:kojleont:teorem_1(0)}

\|u(t)\|_{L_k(\Omega_{r})}\leq \mathcal{M}

\exp\left(-\kappa\Bigl[\frac{r^{p_sm_s}}{

t}\Bigr]^{1/(p_sm_s-1)}\right)\|\varphi\|_{L_k(\Omega)}.

\end{equation}

\end{theorem}



\smallskip



На основе неравенства \eqref{eq:kojleont:teorem_1(0)}

устанавливается оценка снизу скорости убывания решения задачи

\eqref{eq:kojleont:ur}--\eqref{eq:kojleont:gu} при $t\to \infty$.





\smallskip



\begin{theorem}[2]\label{eq:kojleont:the_0} Пусть выполнено условие  \eqref{eq:kojleont:supp},

тогда существует положительное число $C(\varphi, k, p_1, \widehat{a}, \widehat{b})$ такое$,$ что ограниченное решение

$u(t,\mathbf{x})$ задачи

\eqref{eq:kojleont:ur}--\eqref{eq:kojleont:gu} при всех $t\geq 0$

подчиняется оценке

\begin{equation}\label{eq:kojleont:teorem_0}

\|u(t)\|_{L_k(\Omega)}\geq

\|\varphi\|_{L_k(\Omega)}\left(C(\varphi)t+1\right)^{-1/(p_1-k)}.

\end{equation}

\end{theorem}



\smallskip



Показано, что наилучшая скорость убывания решений достигается в~сужающихся неограниченных областях,

именно (см. замечание) для решений задачи \eqref{eq:kojleont:ur}--\eqref{eq:kojleont:gu} в~области $\Omega$, удовлетворяющей условию

\begin{equation}

\label{eq:kojleont:mu}

\mu_1=\inf\left\{  \Bigl\|\frac{\partial^{m_1} g}{\partial x_1^{m_1}} \Bigr\|_{L_{p_1}(\Omega)} \mid g({\bf x})\in C_0^{\infty}(\Omega), \|g\|_{L_{k}(\Omega)}=1\right\}>0,\end{equation}

справедлива оценка

\begin{equation}

\label{eq:kojleont:sle_1}

\|u(t)\|_{L_k(\Omega)}\leq{M}t^{-1/(p_{1}-k)},\quad t>0.

\end{equation}



\smallskip



\Section{Вспомогательные утверждения} Пусть $\| {}\cdot{}\|_{p,Q}$ "--- норма в $L_p(Q)$, $p\geq 1$, причём значение $Q=\Omega$ опускается. 

Через $D_a^b=(a, b)\times\Omega$ обозначим цилиндр, значения $a=0$ и $b=\infty$ могут отсутствовать.



Банахово пространство $\stackrel{\circ}{W}\!{}_{k,{\bf p}}^{\;\;{\bf m}}(\Omega)$ определим как пополнение пространства

$C_0^{\infty}(\Omega)$ по норме

$$

\|u\|_{{W}_{k,{\bf

p}}^{\;\;{\bf m}}(\Omega)}=\sum_{\alpha=1}^n\Bigl\|\frac{\partial^{m_\alpha}u}{\partial x_\alpha^{m_\alpha}}

\Bigr\|_{p_{\alpha}}+\|u \|_k.

$$

Банаховы пространства

 $\stackrel{\circ}{W}\!{_{k,{\bf p}}^{0,{\bf m}}}(D^T)$,

 $\stackrel{\circ}{W}\!{_{k,{\bf p}}^{1,{\bf m}}}(D^T)$

определим как пополнения пространства

$C_0^{\infty}(D_{-1}^{T+1})$, соответственно,  по нормам

\begin{gather*}

 \| u \|_{W{_{k,{\bf p}}^{0,{\bf m}}}(D^T)}=\|u

\|_{k,D^{T}}+\sum_{\alpha=1}^n\Bigl\| \frac{\partial^{m_\alpha}u}{\partial

x_\alpha^{m_\alpha}}\Bigr\|_{p_{\alpha},D^{T}},\\

\| u

\|_{W{_{k,{\bf p}}^{1,{\bf m}}}(D^T)}=\| u \|_{W{_{k,{\bf p}}^{0,{\bf

m}}}(D^T)}+\|u_t \|_{k,D^{T}}.

\end{gather*}



\smallskip



\begin{definition} Обобщённым решением задачи \eqref{eq:kojleont:ur}--\eqref{eq:kojleont:gu} назовём функцию $u(t,{\bf x})$ такую, что  при всех $T>0$ $u(t,{\bf x})\in {\stackrel {\circ} {W}\!^{0,{\bf m}}_{k,{\bf

p}}}(D^T)$, при 

$k\in(1, 2)$ $u_t(t,{\bf x})\in L_k(D^T),$  а~при

$k\geq 2$ $|u|^{k-2}u_t\in L_{k'}(D^T)$, $k'={k}/(k-1)$, и

удовлетворяет интегральному тождеству

\begin{multline}

\label{eq:kojleont:intt}

\int_{D^{T}}

\biggl(-|u|^{k-2}uv_{t}+\sum_{\alpha=1}^{n} a_{\alpha}\left(\Bigl(\frac{\partial^{m_\alpha}u}{\partial x_\alpha^{m_\alpha}}\Bigr)^2\right)\frac{\partial^{m_\alpha}u}{\partial

x_\alpha^{m_\alpha}} \frac{\partial^{m_\alpha}v}{\partial

x_\alpha^{m_\alpha}} \biggr)d{\bf x}dt+

\\

+\int_ \Omega |u(T,{\bf x})|^{k-2}u(T,{\bf x}) v(T,{\bf

x})d{\bf x}=\int_ \Omega |\varphi({\bf

x})|^{k-2}\varphi({\bf x}) v(0,{\bf x})d{\bf x}

\end{multline}

для любой функции $v(t,{\bf x}) \in \stackrel{\circ}{W}{_{k,{\bf

p}}^{1,{\bf m}}}(D^T).$



\end{definition}



\smallskip



\begin{theorem}[3]Пусть $\varphi({\bf x})\in \stackrel{\circ}{W}\!{}_{k,{\bf p}}^{\;\;\;\bf{m}}(\Omega),$ $p_1\geq k,$ $k>1,$ тогда существует обобщённое решение $u(t,{\bf x})$ задачи \eqref{eq:kojleont:ur}--\eqref{eq:kojleont:gu}$.$ 

При этом справедливы неравенства

\begin{gather}

\label{eq:kojleont:theorem_4_3}

(k-1)\|u(t)\|^k_k+k\overline{a}\sum_{\alpha=1}^n\int_0^t \Bigl\|\frac{\partial^{m_\alpha}u(\tau)}{\partial

x_\alpha^{m_\alpha}}\Bigr\|_{p_{\alpha}}^{p_{\alpha}}d\tau\leq

(k-1)\|\varphi\|^k_k,\quad t\geq 0,\\

\label{eq:kojleont:theorem_4_4}

(k-1)\frac{d}{dt}\|u(t)\|^k_k+k\overline{a}\sum_{\alpha=1}^n

\Bigl\|\frac{\partial^{m_\alpha}u(t)}{\partial x_\alpha^{m_\alpha}}\Bigr\|_{p_{\alpha}}^{p_{\alpha}}\leq 0,\quad

t>0.

\end{gather}

\end{theorem}



\smallskip



Решение задачи \eqref{eq:kojleont:ur}--\eqref{eq:kojleont:gu}

строится методом галёркинских приближений, который ранее был предложен Ф.~Х.~Мукминовым, Э.~Р.~Андрияновой (см.~\cite{kojleont:bib:and_muk}) для модельного изотропного параболического уравнения с~двойной нелинейностью и~обобщён

авторами статьи на некоторый класс анизотропных уравнений (см.~\cite{kojleont:bib:umg_2011,kojleont:bib:umg_2012}

).



\smallskip



\begin{newthm}{Предложение 1} Для любой функции $g(x)\in L_p(\mathbb{R}_{1}^+),$ $D^mg(x)\in L_p(\mathbb{R}_{1}^+)$

справедливо неравенство

\begin{equation}

\label{eq:kojleont:N-G} \|D^jg\|_{p,\mathbb{R}_{1}^+}\leq C

\|D^mg\|_{p,\mathbb{R}_{1}^+}^{\frac{j}{m}}

\|g\|_{p,\mathbb{R}_{1}^+}^{1-\frac{j}{m}},\quad j=0, 1, \dots, m,

\end{equation}

где постоянная $C$ зависит от $p,$ $m,$ $j.$

\end{newthm}



\smallskip



\begin{proof}

Для произвольной функций $h(x)\in L_p(\mathbb{R}_{1})$, $D^m h\in

L_p(\mathbb{R}_{1})$  справедливо одномерное неравенство

Ниренберга"--~Гальярдо~\cite{kojleont:bib:nirenberg}, которое

запишем в частном случае:

$$

\|D^jh\|_{p,\mathbb{R}_{1}}\leq c_1

\|D^mh\|_{p,\mathbb{R}_{1}}^{\frac{j}{m}}

\|h\|_{p,\mathbb{R}_1}^{1-\frac{j}{m}},\quad

j=0, 1, \dots, m.

$$

Пусть функция $h(x)$ "--- продолжение функции $g(x)$ такое, что

$$

\|D^jh\|_{p,\mathbb{R}_{1}}\leq c_2\|D^j g\|_{p,\mathbb{R}_{1}^+},\quad

j=\overline{0,m}.

$$

Соединяя эти неравенства, получаем неравенство

\eqref{eq:kojleont:N-G}.

\end{proof}



\smallskip



\Section{ Доказательство теорем}  Будем предполагать, что решение

задачи \eqref{eq:kojleont:ur}--\eqref{eq:kojleont:gu} ограничено,

а именно, справедливо неравенство

\begin{equation}

\label{eq:kojleont:prop}

{\rm vrai}\sup\limits_{D}\mid u(t,{\bf x})\mid\leq B<\infty.

\end{equation}



\begin{proof}[Доказательство \ теоремы \ 1]

Пусть $\theta(x)$, $x>0$ "--- гладкая неотрицательная функция, равная единице при $x\ge 1$, нулю при $x\le 0$. 

Тогда для функции $\xi(x_s)=\theta((x_s-r)/\rho)$ справедливы соотношения 

\begin{equation}

\label{eq:kojleont:the_1(0)} 

\Bigl|\frac{d^{j}\xi}{d x_s^{j}}\Bigr|\leq C_1/\rho^j,\quad x\in(r,r+\rho),\quad

\frac{d^{j}\xi}{d x_s^{j}}= 0,\quad x\not\in(r,r+\rho),\quad j \in

\mathbb{N}.

\end{equation}



Для $k\in(1, 2)$ в \eqref{eq:kojleont:intt} можно взять $v=u\xi$, тогда, пользуясь равенством

$$\int_{\Omega}|u|^{k-2}uu_{\tau}d{\bf x}=\frac{1}{k}\frac{d}{d\tau}\biggl(\int_{\Omega}|u|^kd{\bf x}\biggr),$$

получим  тождество

\begin{equation}

\frac{k-1}{k}\int_{ \Omega}|u|^{k} \xi \Bigl |_{ \tau=0}^{  \tau=t}d{\bf x}

  +

  \sum\limits_{\alpha=1}^n\int_{D^{t}}

  a_{\alpha}\left(\Bigl(\frac{\partial^{m_\alpha}u}{\partial x_\alpha^{m_\alpha}}\Bigr)^2\right)

  \frac{\partial^{m_\alpha}u}{\partial x_\alpha^{m_\alpha}}

\frac{\partial^{m_\alpha}(u\xi)}{\partial x_\alpha^{m_\alpha}} d{\bf x}d\tau =0.

\label{eq:kojleont:teorem_3_2(14_1)}

\end{equation}

 Для $k\geq 2$ можно выполнить интегрирование по частям в первом интеграле тождества \eqref{eq:kojleont:intt},

 в результате получим равенство

$$

\int_{D^{\tau}} \biggl(

(|u|^{k-2}u)_{\tau}v+

\sum_{\alpha=1}^{n} 

a_{\alpha}\left(\Bigl(\frac{\partial^{m_\alpha}u}{\partial x_\alpha^{m_\alpha}}\Bigr)^2\right)

\frac{\partial^{m_\alpha}u}{\partial x_\alpha^{m_\alpha}} \frac{\partial^{m_\alpha}v}{\partial x_\alpha^{m_\alpha}}

\biggr)d{\bf x}d\tau=0.

$$

Положим $v=u\xi$, тогда, ввиду справедливости равенства

$$

\int_{\Omega}(|u|^{k-2}u)_{\tau}ud{\bf x}=\frac{k-1}{k}\frac{d}{d\tau}\biggl(\int_{\Omega}|u|^kd{\bf x}\biggr),

$$ выводим  соотношение \eqref{eq:kojleont:teorem_3_2(14_1)}.



Далее, применяя \eqref{eq:kojleont:us1}, из

\eqref{eq:kojleont:teorem_3_2(14_1)} получаем (с учётом того, что $\xi\varphi=0$)

\begin{multline}

\label{eq:kojleont:the_1(04)} 

\frac{k-1}{k}\int_{\Omega}|u(t,\mathbf{x})|^k \xi(x_s) d\mathbf{x} +

\overline{a}\sum\limits_{\alpha=1}^n \int_{D^{t}} \xi

\Bigl|\frac{\partial^{m_\alpha}u}{\partial x_\alpha^{m_\alpha}}\Bigr|^{p_{\alpha}}d\mathbf{x}d \tau\leq \\

\leq \widehat{a}\int_{D^{t}}\sum\limits_{j=1}^{m_s}

\Bigl|\frac{\partial^{m_s}u}{\partial x_s^{m_s}}\Bigr|^{p_{s}-1}C_{m_s}^j \Bigl|\frac{\partial^{m_s-j}u}{\partial x_s^{m_s-j}}\Bigr| \Bigl|\frac{d^{j}\xi}{d x_s^{j}}\Bigr| d\mathbf{x}d \tau\equiv

 I^t.

\end{multline}

Используя \eqref{eq:kojleont:the_1(0)},  оценим интеграл

$$

 I^t\leq\sum\limits_{j=1}^{m_s}\frac{C_2}{\rho^j}\int_{0}^{t}

   \int_{ \Omega_r^{r+\rho}}

   \Bigl|\frac{\partial^{m_s}u}{\partial x_s^{m_s}}\Bigr|^{p_s-1}

   \Bigl|\frac{\partial^{m_s-j}u}{\partial x_s^{m_s-j}}\Bigr| d\mathbf{x}d

   \tau.

$$



Пусть $\varepsilon>0$ "--- произвольное число такое, что $\varepsilon\rho\geq1$. Тогда, последовательно применяя неравенства Гёльдера, неравенство \eqref{eq:kojleont:N-G} и  Юнга, получаем соотношения

\begin{multline*}

I^t\leq\sum\limits_{j=1}^{m_s}\frac{C_2}{\rho^j} \int _{0}^{t}

 \int _{\mathbb{R}_{n-1}

}\Bigl\|\frac{\partial^{m_s}u}{\partial x_s^{m_s}}\Bigr\|^{p_s-1}_{p_s,(r,r+\rho)}

\Bigl\|\frac{\partial^{m_s-j}u}{\partial x_s^{m_s-j}}\Bigr\|_{p_s,(r,r+\rho)} d\mathbf{x}'_sd \tau\leq \\

\leq\sum\limits_{j=1}^{m_s}\frac{C_3}{\rho^j\varepsilon^j} \int _{0}^{t}

 \int _{\mathbb{R}_{n-1} }\Bigl\|\frac{\partial^{m_s}u}{\partial x_s^{m_s}}\Bigr\|^{p_s-j/m_s}_{p_s,(r,\infty)}\|u\|^{j/m_s}_{p_s,(r,\infty)}

\varepsilon^jd\mathbf{x}'_sd \tau\leq\\

\leq\sum\limits_{j=1}^{m_s}\frac{C_3}{\rho^j\varepsilon^j} \int _{0}^{t}

 \int _{\mathbb{R}_{n-1}

}\left[\frac{m_sp_s-j}{m_sp_s}\Bigl\|\frac{\partial^{m_s}u}{\partial x_s^{m_s}}\Bigr\|^{p_s}_{p_s,(r,\infty)}+\varepsilon^{m_sp_s}\frac{j}{m_sp_s}\|u\|^{p_s}_{p_s,(r,\infty)}

\right] d\mathbf{x}'_sd \tau\leq \\

\leq\frac{C_4}{\rho\varepsilon} \int _{0}^{t}

 \int _{\Omega_r}\left[\Bigl|\frac{\partial^{m_s}u}{\partial x_s^{m_s}}\Bigr|^{p_s}+\varepsilon^{m_sp_s}|u|^{p_s}\right]d\mathbf{x}d

\tau, \end{multline*} из которых с учётом \eqref{eq:kojleont:prop} имеем

\begin{equation}

\label{eq:kojleont:the_1(05)}

I^t\leq\frac{C_4}{\rho\varepsilon}\left( \int _{0}^{t}

 \int _{\Omega_r}\Bigl|\frac{\partial^{m_s}u}{\partial x_s^{m_s}}\Bigr|^{p_s}d\mathbf{x}d

\tau+\varepsilon^{m_sp_s}B^{p_s-k} \int _{0}^{t}

 \int _{\Omega_r}|u|^{k}d\mathbf{x}d \tau\right).

\end{equation}

Соединяя \eqref{eq:kojleont:the_1(04)}, \eqref{eq:kojleont:the_1(05)}, выводим неравенство

\begin{multline}

\label{eq:kojleont:the_1(06)} \frac{k-1}{k}\|u(t)\|^{k}_{k,

\Omega_{r+\rho}}  + \overline{a}\sum\limits_{\alpha=1}^n

 \int _0^t \Bigl\|\frac{\partial^{m_\alpha}u}{\partial x_\alpha^{m_\alpha}}\Bigr\|^{p_{\alpha}}_{p_{\alpha},\Omega_{r+\rho}}d

\tau \leq \\

\leq\frac{C_5}{\rho\varepsilon}\left( \int _{0}^{t} \Bigl\|\frac{\partial^{m_s}u}{\partial x_s^{m_s}}\Bigr\|^{p_s}_{p_s,\Omega_r}d

\tau+\varepsilon^{m_sp_s} \int _{0}^{t}

\|u\|^{k}_{k,\Omega_r}d \tau\right).

\end{multline}



Введём обозначение

$$

F_r(t)=\|u(t)\|^{k}_{k, \Omega_{r}} + \sum\limits_{\alpha=1}^n

 \int _0^t \Bigl\|\frac{\partial^{m_\alpha}u}{\partial x_\alpha^{m_\alpha}}\Bigr\|^{p_{\alpha}}_{p_{\alpha},\Omega_{r}}d

\tau,

$$

тогда \eqref{eq:kojleont:the_1(06)} можно переписать в виде

\begin{equation}

\label{eq:kojleont:the_1(07)}

F_{r+\rho}(t)\leq

\frac{C_6}{\varepsilon\rho}\left(F_r(t)+\varepsilon^{p_sm_s} \int _0^tF_r(\tau)d\tau\right).

\end{equation}



Отметим, что если $\varepsilon\rho<1$, то неравенство

\eqref{eq:kojleont:the_1(07)} также выполняется, так как в этом

случае

$$

F_{r+\rho}(t)\leq F_r(t)\leq\frac{1}{\varepsilon\rho}

F_r\leq\frac{C_6}{\varepsilon\rho}\left(F_r(t)+\varepsilon^{p_sm_s} \int _0^tF_r(\tau)d\tau\right).

$$

Таким образом, \eqref{eq:kojleont:the_1(07)} доказано для

произвольного $\varepsilon>0$.



Далее индукцией по $l=0, 1, \dots $ установим неравенство

\begin{equation}

\label{eq:kojleont:the_1(08)} F_{R_0+l\rho}(t)\leq C_7 

\Bigl(\frac{2C_6}{\rho}\Bigr)^lt^{l/(p_sm_s)}\biggl\{\prod\limits_{i=0}^{l-1}

\left(1+{i}/{(p_sm_s)}\right)\biggr\}^{-1/(p_sm_s)}  \|\varphi\|_k^k.

\end{equation}

В качестве нулевого шага индукции из неравенства

\eqref{eq:kojleont:theorem_4_3} для любого $t>0$ имеем неравенство

$F_{R_0}(t)\leq C_7\|\varphi\|_k^k$. Предположим, что

\eqref{eq:kojleont:the_1(08)} справедливо для некоторого целого

$l\geq 0$. Подставляя в \eqref{eq:kojleont:the_1(07)}

$$\varepsilon=\left[\frac{1+{l}/(p_sm_s)}{t}\right]^{1/(p_sm_s)},\quad r=R_0+l\rho,$$ с учётом \eqref{eq:kojleont:the_1(08)} получаем

\begin{multline*}

F_{R_0+(l+1)\rho}(t)\leq

C_72^l\Bigl(\frac{C_6}{\rho}\Bigr)^{l+1}t^{1/(p_sm_s)}

\biggl\{\prod\limits_{i=0}^{l}\left(1+{i}/(p_sm_s)\right)\biggr\}^{-1/(p_sm_s)}

\|\varphi\|_k^k\times \\

\times\left\{t^{l/(p_sm_s)}+\frac{1+l/(p_sm_s)}{t} \int _0^t\tau^{l/(p_sm_s)}d\tau\right\}=\\

=C_7\Bigl(\frac{2C_6} {\rho}\Bigr)^{l+1}t^{(l+1)/(p_sm_s)}

\biggl\{\prod\limits_{i=0}^{l}\left(1+{i}/(p_sm_s)\right)\biggr\}^{-1/(p_sm_s)}

\|\varphi\|_k^k.

\end{multline*}

Неравенство \eqref{eq:kojleont:the_1(08)} доказано.



Положим $\rho=(r-R_0)/l$. Используя неравенство Стирлинга, из \eqref{eq:kojleont:the_1(08)} нетрудно получить

\begin{equation}

\label{eq:kojleont:the_1(09)}

F_{r}(t)\leq C_9\exp\Bigl(-\frac{l}{p_sm_s}\ln

\frac{(r-R_0)^{p_sm_s}}{C_8 tl^{p_sm_s-1}}\Bigr)\|\varphi\|_k^k.

\end{equation}

Полагая $l$ равным целой части выражения $$\displaystyle\Bigl[\frac{(r-R_0)^{p_sm_s}}{eC_8

t}\Bigr]^{1/(p_sm_s-1)},$$ из неравенства \eqref{eq:kojleont:the_1(09)} получим

\begin{equation}

\label{eq:kojleont:the_1(010)}

F_{r}(t)\leq C_{11}\exp\left(-C_{10}\Bigl[\frac{(r-R_0)^{p_sm_s}}{t}\Bigr]^{1/(p_sm_s-1)}\right)\|\varphi\|_k^k.

\end{equation}

В итоге при $r\geq 2R_0$ из \eqref{eq:kojleont:the_1(010)} следует оценка \eqref{eq:kojleont:teorem_1(0)}.

\end{proof}



\smallskip



\textit{Д\,о\,к\,а\,з\,а\,т\,е\,л\,ь\,с\,т\,в\,о \ т\,е\,о\,р\,е\,м\,ы~2}\/ осуществляется аналогично доказательству

для случая уравнения второго порядка (подробнее

см.~\cite{kojleont:bib:umg_2011,kojleont:bib:umg_2012}).



\smallskip



\begin{newthm}{Замечание} Если выполнено условие \eqref{eq:kojleont:mu}$,$  то для  решения $u(t,{\bf x})$ задачи \eqref{eq:kojleont:ur}--\eqref{eq:kojleont:gu} справедлива оценка  \eqref{eq:kojleont:sle_1}.

\end{newthm}



\smallskip



\begin{proof} Из \eqref{eq:kojleont:theorem_4_4}  следует, что

$$

\frac{d}{dt}\|u(t)\|^k_k\leq

-\frac{\overline{a}k}{k-1}\sum_{\alpha=1}^n 

\Bigl\|\frac{\partial^{m_\alpha}u(t)}{\partial x_\alpha^{m_\alpha}} \Bigr\|_{p_{\alpha}}^{p_{\alpha}}\leq-

\frac{\overline{a}k}{k-1}\Bigl\|\frac{\partial^{m_1}u(t)}{\partial x_1^{m_1}}\Bigr\|_{p_1}^{p_1}\leq -\frac{\overline{a}k}{k-1}\mu_1^{p_1}\|

u\|^{p_1}_k.

$$

Решая это дифференциальное неравенство, получим оценку

$$

\|u(t)\|_k\leq

t^{-1/(p_1-k)}\Bigl(\frac{(p_1-k)\overline{a}}{k-1}\mu_1^{p_1}\Bigr)^{-1/(p_1-k)},\quad t>0,

$$

из которой следует неравенство \eqref{eq:kojleont:sle_1}.

\end{proof}





\thanks{Работа выполнена при финансовой поддержке РФФИ (проект \No\ 13--01--0081-а).}



\begin{thebibliography}{9}



\RBibitem{kojleont:bib:kojevnik2000}

\by Л.~М.~Кожевникова, Ф.~Х.~Мукминов

\paper Оценки скорости стабилизации при \mbox{$t\to\infty$} решения первой смешанной задачи для квазилинейной системы параболических уравнений второго порядка

\jour Матем. сб.

\yr 2000

\vol 191

\issue 2

\pages 91--131

\mathnet{http://mi.mathnet.ru/msb454}

\crossref{http://dx.doi.org/10.4213/sm454}

\mathscinet{http://www.ams.org/mathscinet-getitem?mr=1751776}

\zmath{http://www.zentralblatt-math.org/zmath/search/?an=Zbl 0957.35068}

\transl англ. пер.:~

\by L.~M.~Kozhevnikova, F.~Kh.~Mukminov

\paper Estimates of the stabilization rate as $t\to\infty$ of solutions of the~first mixed problem for a~quasilinear system of second-order parabolic equations

\jour Sb. Math.

\yr 2000

\vol 191

\issue 2

\pages 235--273

\crossref{http://dx.doi.org/10.1070/sm2000v191n02ABEH000454}

\isi{http://gateway.isiknowledge.com/gateway/Gateway.cgi?GWVersion=2&SrcApp=PARTNER_APP&SrcAuth=LinksAMR&DestLinkType=FullRecord&DestApp=ALL_WOS&KeyUT=000087494000009}





\RBibitem{kojleont:bib:kojevnik2005}

\by Л.~М.~Кожевникова

\paper Стабилизация решения первой смешанной задачи для~эволюционного квазиэллиптического уравнения

\jour Матем. сб.

\yr 2005

\vol 196

\issue 7

\pages 67--100

\mathnet{http://mi.mathnet.ru/msb1377}

\crossref{http://dx.doi.org/10.4213/sm1377}

\mathscinet{http://www.ams.org/mathscinet-getitem?mr=2188370}

\zmath{http://www.zentralblatt-math.org/zmath/search/?an=Zbl 1085.35050}

\transl англ. пер.:~

\by L.~M.~Kozhevnikova

\paper Stabilization of a~solution of the first mixed problem for a~quasi-elliptic evolution equation

\jour Sb. Math.

\yr 2005

\vol 196

\issue 7

\pages 999--1032

\crossref{http://dx.doi.org/10.1070/SM2005v196n07ABEH000946}

\isi{http://gateway.isiknowledge.com/gateway/Gateway.cgi?GWVersion=2&SrcApp=PARTNER_APP&SrcAuth=LinksAMR&DestLinkType=FullRecord&DestApp=ALL_WOS&KeyUT=000232881000004}







\RBibitem{kojleont:bib:kar_koj}

\by Р.~Х.~Каримов, Л.~М.~Кожевникова

\paper Стабилизация решений квазилинейных параболических уравнений второго порядка в~областях с~некомпактными границами

\jour Матем. сб.

\yr 2010

\vol 201

\issue 9

\pages 3--26

\mathnet{http://mi.mathnet.ru/msb7602}

\crossref{http://dx.doi.org/10.4213/sm7602}

\mathscinet{http://www.ams.org/mathscinet-getitem?mr=2760458}

\adsnasa{http://adsabs.harvard.edu/cgi-bin/bib_query?2010SbMat.201.1249K}

\transl англ. пер.:~

\by R.~Kh.~Karimov, L.~M.~Kozhevnikova

\paper Stabilization of solutions of quasilinear second order parabolic equations in domains with non-compact boundaries

\jour Sb. Math.

\yr 2010

\vol 201

\issue 9

\pages 1249--1271

\crossref{http://dx.doi.org/10.1070/SM2010v201n09ABEH004111}

\isi{http://gateway.isiknowledge.com/gateway/Gateway.cgi?GWVersion=2&SrcApp=PARTNER_APP&SrcAuth=LinksAMR&DestLinkType=FullRecord&DestApp=ALL_WOS&KeyUT=000285190100001}







\RBibitem{kojleont:bib:and_muk}

\by Э.~Р.~Андриянова, Ф.~Х.~Мукминов

\paper Оценка снизу скорости убывания решения параболического уравнения с~двойной нелинейностью

\jour Уфимск. матем. журн.

\yr 2011

\vol 3

\issue 3

\pages 3--14

\mathnet{http://mi.mathnet.ru/ufa99}

\misctransl

\by E.~R.~Andriyanova, F.~Kh.~Mukminov

\paper The lower estimate of decay rate of solutions for doubly nonlinear parabolic equations

\jour Ufimsk. Mat. Zh.

\yr 2011

\vol 3

\issue 3

\pages 3--14









\RBibitem{kojleont:bib:umg_2011}

\by Л.~М.~Кожевникова, А.~А.~Леонтьев

\paper Оценки решения анизотропного параболического уравнения с~двойной нелинейностью

\jour Уфимск. матем. журн.

\yr 2011

\vol 3

\issue 4

\pages 64--85

\mathnet{http://mi.mathnet.ru/ufa119}

\misctransl

\by L.~M.~Kozhevnikova, A.~A.~Leontiev

\paper Estimates of solutions of an anisotropic doubly nonlinear parabolic equation

\jour Ufimsk. Mat. Zh.

\yr 2011

\vol 3

\issue 4

\pages 64--85





\RBibitem{kojleont:bib:umg_2012}

\by Л.~М.~Кожевникова, А.~А.~Леонтьев 

\paper Убывание решения анизотропного параболического уравнения с двойной нелинейностью в неограниченных областях

\jour Уфимск. матем. журн.

\yr 2013

\vol 5

\issue 1

\pages 65--83

\misctransl

\paper Decay of solution of anisotropic doubly nonlinear parabolic equation in unbounded domains

\by L.~M.~Kozhevnikova, A.~A.~Leontiev 

\jour Ufimsk. Mat. Zh.

\yr 2013

\vol 5

\issue 1

\pages 65--83









\Bibitem{kojleont:bib:nirenberg}

\by L.~Nirenberg

\paper On elliptic partial differential equations

\jour Ann. Sc. Norm. Super. Pisa, Sci. Fis. Mat., III. Ser

\yr 1959

\vol 13

\pages 115--162





\end{thebibliography}





\noindent

\hskip6.5cm \thedate \par\noindent

\hskip6.5cm\therevdate

%





\newpage





\chead[\fancyplain{}{\scriptsize \sl L.~M.~K\,o\,z\,h\,e\,v\,n\,i\,k\,o\,v\,a,

A.~A.~L\,e\,o\,n\,t\,i\,e\,v }]{\fancyplain{}{\scriptsize \sl  Solutions of anisotropic parabolic equations with double non-linearity \dots}}



\makeengtitlem







 \newpage

%

% \tableofengcontents

%

%\end{document}

